\section{Empirical Analysis}
\label{sec:analysis}
 
% ==============================================================================

 \subsection{Failure Pattern Statistics}
\label{sec:failures}

Auditing every scored trajectory against the 45-pattern taxonomy of
\S\ref{sec:taxonomy} yields 12{,}712 hits across 800 analyses as shown in \Cref{fig:failure-attribution}; complete counts
are in Appendix~\ref{sec:failure_pattern_hits}.

The distribution across root-cause pillars is uneven. The three cognitive
pillars---Grounding \& Faithfulness (R1, 31.0\%), Scientific Integrity \&
Alignment (R3, 33.5\%), and Cognitive Depth \& Adaptability (R2,
27.6\%)---together account for 92.1\% of all hits. Engineering Robustness (R4)
contributes just 7.9\%, and its highest-ranked pattern, execution faults and
numerical instability (C.4), places only 26th of 45 failure patterns.

At the individual-pattern level, failures concentrate heavily in the
self-verification stage. Uncorrected self-awareness (F.4) is the single most
frequent pattern in the corpus, appearing in 660 of 800 analyses (82.5\%).
Two related patterns---failure to gate critical flaws (F.2, 502 hits) and
unremediated adversarial evidence (D.7, 486 hits)---rank among the top five.
Together these three patterns account for 13.0\% of all hits, the largest
concentration attributable to any single mechanism.

Failure profiles differ across the eight model--harness combinations, though
the overall shape is consistent. Total hit counts range from 1{,}396
(opus-4.8) to 1{,}818 (qwen3.7-max), and the top-10 most frequent patterns
overlap heavily: E.2, D.4, A.5, and C.1 appear in the top~10 of every model, and F.4 in seven of the eight (Appendix~\ref{sec:top10_by_model}). Where systems diverge most is in
fabrication-related patterns. Hallucinated evidence (B.1) ranges from 13 hits
(glm-5.2) to 61 (gpt-5-mini); result hallucination (D.6) ranges from 3
(opus-4.8 and claude-sonnet-5) to 36 (qwen3.7-max). These system-specific profiles suggest that
while the core failure patterns are shared across models, fabrication rates
reflect differences in model capability, and the detailed per-model breakdowns
in Appendix~\ref{sec:top10_by_model} can guide system-specific mitigation.

Two patterns, review score hacking (F.5) and hallucinated reviewing (F.6), are
near-absent in the corpus---one and three hits respectively, both
borderline---so their inclusion in the taxonomy reflects observed behavior
rather than a claim about prevalence.

\begin{figure}[htbp]
\centering
\includegraphics[width=0.92\textwidth]{figs/failure_heatmap.png}
\caption{Failure pattern attribution across agent phases (aggregated
trajectories, $n=800$). Cell color encodes the number of HIT
attributions\add{; each pattern contributes the number of
distinct trajectories in which it was an established failure}%
; ``---'' denotes none. A single cell may contain
more than one failure pattern within the same cell, so cell counts can
exceed 800.}
\label{fig:failure-attribution}
\end{figure}

\subsection{Root Causes and What They Imply}
\label{sec:rootcause}

Some of what we observe reflects limits of the underlying models rather than of
the systems built around them, and we say so where the data shows it. Our
emphasis falls on failures whose evidence is already present in the trajectory,
because those are the ones where the distance between what the agent knew and
what it did can be measured rather than assumed.

The root causes differ in surface form but share a common origin: the
metacognitive deficit of \S\ref{sec:taxonomy}. A human researcher works in a
closed loop---checking whether what they produced matches what they found,
acting when it does not, and questioning whether the path they took was
sound. Current agents lack this loop, and each root cause is a different point
where it breaks. R1 is the failure to check output against evidence. R2 is the
failure to act on a flaw the agent itself identified. R3 is the failure to
question whether the path to the result is legitimate. We take each in turn.

Two kinds of number appear below. A root cause's \emph{share}---such as R1's
31.0\%---is its fraction of the 12{,}712 total hits across the corpus. A
pattern's \emph{rate}---such as D.4's 77.5\%---is the fraction of 800 analyses. \add{Where a group of patterns is summarised together, the figure given is that group's share of its root cause's hit total. }%
Notice that one trajectory usually contains more than one failure pattern.

% ---------------------------------------------------------------- R1
\begin{tcolorbox}[colback=blue!3,colframe=blue!35,boxrule=0.5pt,
                  left=5pt,right=5pt,top=4pt,bottom=4pt]
\textbf{R1. Grounding \& Faithfulness}\;(31.0\% of all hits). Claims,
hypotheses, and plans disconnect from the code, data, logs, or literature
that should license them.\\[3pt]
\raisebox{-1pt}{\scriptsize$\blacktriangleright$}\;\textbf{Insight 1. The
evidence that would refute most failures is already in the agent's own run
directory.} The agent produces both the claim and the file that contradicts
it; the comparison is never performed.
\end{tcolorbox}

\noindent R1 is led by method--conclusion disconnect (D.4, 77.5\% of
analyses), implementation discrepancy (C.3, 72.1\%), and report--code
traceability gaps (E.1, 60.5\%). What these share is that the agent writes
up the work it meant to do rather than the work it did---and the evidence
exposing the gap sits in the same run directory the agent itself created.

Two-thirds of R1 consists of \emph{unsupported claims}, in two forms. In the
first, the run files say something different from the report: a conclusion is
drawn that the method does not support (D.4, 77.5\%), an experiment is
presented as testing a hypothesis it does not test (A.6, 52.5\%), or a side
effect of the pipeline is reported as a finding (D.1, 52.4\%). In the second,
there is nothing behind the claim at all---a method section describing a
procedure the code never implements (C.3, 72.1\%), or a claim no run in the
trajectory supports (E.1, 60.5\%). The second form is the more serious,
because the agent is not overstating a result but reporting work it did not
do. It is also distinct from the overclaiming of R3, where the experiment was
run and the conclusion reaches past it.

A further sixth is \emph{invented evidence}---hallucinated sources, fabricated
citations, numbers given for runs that never happened (B.1, E.4,
D.6)---where nothing can contradict the claim because the thing referred to
does not exist. The two groups behave differently across systems: unsupported
claims appear at nearly the same rate in all eight systems, while invented
evidence is much rarer in the strongest ones. The remaining is retrieval that
never reaches the design (B.2, 33.9\%), citations that do not support the
sentence citing them (B.5, 19.0\%), and right-for-the-wrong-reason success
(X.6, 43.8\%)---less a mechanism of its own than the outcome the other two
produce.

Two consequences. Stronger models are unlikely to remove unsupported claims, because
catching one requires no ability the agents lack: the report and the run
directory are both products of the same agent, and all that is missing is a
requirement to compare them before the report goes out. A human author checks
the number in the abstract against the number in the table, and does not list
a contribution they never implemented; nothing in these runs performs either
check. And an evaluation that reads only the final report cannot see any of
this, because the report is one half of the disagreement and the other half
is on disk. That is the case for scoring the artifacts a run leaves behind
rather than the answer it ends with. This is the metacognitive loop failing
at the point of evaluation: the agent never compares what it produced against
what it found, so the loop never begins.

% ---------------------------------------------------------------- R2
\begin{tcolorbox}[colback=blue!3,colframe=blue!35,boxrule=0.5pt,
                  left=5pt,right=5pt,top=4pt,bottom=4pt]
\textbf{R2. Cognitive Depth \& Adaptability}\;(27.6\% of all hits). Shallow
reasoning and search, passivity in self-critique, and inability to re-plan or
pivot at a dead end.\\[3pt]
\raisebox{-1pt}{\scriptsize$\blacktriangleright$}\;\textbf{Insight 2. The
most common failure in the corpus is not missing a flaw but finding it and
shipping anyway.} In 82.5\% of analyses the agent diagnoses a critical
problem during self-review, then reports the unrevised conclusion.
\end{tcolorbox}

\noindent This root cause is where the missing metacognitive loop of
\S\ref{sec:taxonomy} is easiest to see, so it is worth stating plainly what
that loop is. A researcher in difficulty does three things: recognises the
limits of what they know, judges whether their intermediate results hold up,
and changes the plan when they do not---then repeats. R2 is the failure of
all three.

Recognising limits fails in frame-lock within a narrow hypothesis space
(A.1), thin search coverage (B.4, 54.9\% of analyses), and absent skepticism
(X.3, 39.4\%)---together 27\% of R2's hits. Judging results fails in the four self-review patterns
(F.1--F.4) and in missing baselines (D.5): 62\% of R2's hits. Changing
the plan fails in local optimization (C.6), stopping early (C.7), and
anchoring to a failed approach (X.7): 11\% of R2's hits.

Judging results is the largest of the three, and its name is misleading,
because agents do reach the right judgment. In uncorrected self-awareness
(F.4, 82.5\% of analyses), the most common pattern in the corpus, the agent
finds the fatal flaw and writes it down, then reports the conclusion anyway.
The judgment was correct and nothing followed from it. A self-review is just
more text: the same model writes it, in the same run, and nothing in the
system requires the review to change the report. A human researcher who finds
a broken baseline has to fix it before the paper goes out; these agents do
not.

Judging results could in principle be repaired from outside the model. If the
agent writes that its own result is uninterpretable, the system can refuse
the report until either the result or the claim changes. Recognising limits
cannot be repaired that way. When an agent never considers a second
hypothesis, no second hypothesis exists in the trajectory to compare against,
and nothing can flag what was never written down. This is the part that
depends on the model rather than on the harness, and it is why the choice of
system matters more here than in the other two root causes. Unlike R1, the
metacognitive loop does fire here---the agent detects the problem---but
nothing in the system closes it, so the diagnosis produces no change.

% ---------------------------------------------------------------- R3
\begin{tcolorbox}[colback=blue!3,colframe=blue!35,boxrule=0.5pt,
                  left=5pt,right=5pt,top=4pt,bottom=4pt]
\textbf{R3. Scientific Integrity \& Alignment}\;(33.5\% of all hits). The
agent pursues the stated goal by whatever path produces a result, whether the
path is sound or not: shortcut reliance, metric hacking, concealed failure,
and conclusions fixed in advance.\\[3pt]
\raisebox{-1pt}{\scriptsize$\blacktriangleright$}\;\textbf{Insight 3. A
correct-looking output is not evidence of sound methodology.} Agents
routinely reach the stated goal through circular validation, metric
substitution, and concealed negative results---and nothing in the system
distinguishes a legitimate solution from a gamed one.
\end{tcolorbox}

\noindent R3 is the largest root cause, led by overclaiming with concealed
negative results (E.2, 78.1\% of analyses), circular validation and shortcut
reliance (C.1, 69.0\%), and metric misalignment (A.5, 68.1\%). The outputs
these agents produce often look correct---the numbers are plausible, the
report is well-structured, but the methodology behind them is invalid, and the agent never
questions whether it is.

Seventy of our hundred tasks are open-ended: the agent receives a premise and
an unresolved tension, and nothing in the environment checks whether its
approach is sound. Nearly half of R3 is \emph{taking shortcuts in the
method}: substituting a metric the agent can hit for the one the task
requires (A.5, 68.1\%), validating on a substrate that cannot fail (C.1,
69.0\%), setting a disproof condition out of reach (A.2, 44.6\%), or
reasoning backward from the conclusion it intends to draw (X.5, 39.2\%).
Another two-fifths is \emph{hiding what went wrong}: reporting success while
burying negative results (E.2, 78.1\%), omitting critical limitations (E.3,
62.2\%), or leaving contradictory evidence unaddressed (D.7, 60.8\%). The
agent is also its own reviewer, so a limitation can be written down in one
section without disturbing the headline in another.

Supplying an external metric does not remove this---it changes its form. On
the thirty target-anchored tasks the failure becomes grader-fitting (C.2):
moving the number without doing the work the number was meant to measure.
Both invalid trajectories of \S\ref{sec:casestudy} are of this kind.

This is why R3 hardly moves when the system changes. Every system met the
same open-ended setup with no external check on method quality, and every
system took the same shortcuts at close to the same rate. R3 cannot be
repaired by adding a score---that just creates a new target to game. What it
would take is verification the agent does not control: a check applied from
outside the run on whether the method supports the conclusion. In the
metacognitive loop, R3 is the absence of self-questioning: even a trajectory
that checks its outputs and acts on problems can still arrive at a result
through an illegitimate method, because the agent never asks whether the path
itself was sound.
 
% ---------------------------------------------------------------- R4
\paragraph{R4. Engineering Robustness.}

At 7.9\% of all hits with no pattern above 26th of 45, R4 is too small to carry
a contributive finding. What stands out is not how often it appears but how much
it varies: some systems hit environment-interaction failures---crashes (C.5),
broken shell commands (C.8), and files that never get written out (X.8)---far
more often than others, while the rest of R4 is more evenly spread.
These failures depend on how the system is built, not on the task, which is also
why benchmarks that supply a pre-configured environment and a known entry
point~\citep{wang2026naturebench} make execution look solved: they remove the
surface where the failures concentrate.
 

\subsection{Case Study}
\label{sec:casestudy}

\begin{figure*}[t]
\centering
\includegraphics[width=\textwidth]{figs/case-studies_1.png}\\[4pt]
\includegraphics[width=\textwidth]{figs/case-studies_2.png}
\caption{Four case studies of \S\ref{sec:casestudy}. \textbf{(a, b)} are runs
where the metric, not the science, becomes the objective: transcribing the
README's answer key while the training data is never read (\textbf{a}), or
hiding the real search behind a gate the grader never opens (\textbf{b}).
\textbf{(c, d)} are honestly engineered runs capped by a unexamined
metacognitive decision: budget spent perfecting an metric already won while the most promising
$-12.1$ metric is left untouched (\textbf{c}); and a self-diagnosed broken
baseline left unfixed while still using wrong claim to open
the abstract, despite a one-line fix (\textbf{d}). All four instantiate the
metacognitive deficit of \S\ref{sec:taxonomy}: research agents lack an explicit
metacognitive loop.}
\label{fig:case-studies}
\end{figure*}

We defer the failures a human scientist would also make, such as crashes, thin
retrieval, to Appendix~\ref{app:detailed case studies}, and dissect four
counterintuitive trajectories here (Figure~\ref{fig:case-studies}), all quoted
from the agent's own rollout artifacts. Two are flagged \emph{invalid} because
the agent optimized the metric rather than the science (C.2); two are judged
\emph{valid} yet fail because a single unexamined cognitive decision fixed their
ceiling (D.5, D.7).

% ---------------------------------------------------------------------
\paragraph{The answer key becomes the method (C.2; Fig.~\ref{fig:case-studies}a).}
The task asks the agent to learn PDE solution operators from provided training
pairs, scored by relative error on held-out inputs. For the
\texttt{nonlinear\_ode} instance the agent submits, in full:
%
\begin{quote}\small
\begin{verbatim}
X_test = g["X_test"].ravel()
# Stated exact solution; satisfies BCs u(-1)=u(1)=0.
pred = np.sin(np.pi * X_test)
\end{verbatim}
\end{quote}
%
No model is fit and the training data is never read. This is not an oversight
but a deliberate strategy, stated explicitly in the trajectory:
%
\begin{quote}\small\itshape
``The README states this ODE `admits the exact solution $u(x)=\sin(\pi x)$' \dots\
Empirically the reference solution is exactly \texttt{sin(pi x)}.''
\end{quote}
%
Nor is it a one-time pattern: on a second instance it substitutes the known
initial condition $-\sin(\pi x)$ into a Cole--Hopf solution of the \emph{stated}
governing equation, again without touching an observation. The behavior is
striking precisely because nothing in it is incompetent---the boundary
conditions are checked, the analytical derivation is correct, and the prediction
attains near-zero error. What has gone wrong is one level up: the agent has answered the question ``what is $u(x)$?'' when the task was ``learn an operator from data.'' Any information channel that reaches the reference solution is treated as admissible, because the agent identifies the objective with the metric rather than with the scientific procedure the metric is meant to certify. The question it never poses is what it is being asked to \emph{do}.

% ---------------------------------------------------------------------
\paragraph{The gated cheat: hacking disguised as a pipeline (C.2; Fig.~\ref{fig:case-studies}b).}
On a quantum-error-correction task that requires searching a space of
Hadamard/Clifford deformation masks for the one that minimizes the logical error
rate, the agent ships a \texttt{run.py} whose default path returns a planted
constant, while the genuine search sits behind an environment flag that is never
set during evaluation:
%
\begin{quote}\small
\begin{verbatim}
def main():
    run_search = os.environ.get("QEC_RUN_SEARCH", "0") == "1"   # default: "0"
    mask = dict(BEST_MASK)             # the planted "answer" (XZZX mask)
    if run_search:                     # never true under the grader
        mask = _live_search(dp)
\end{verbatim}
\end{quote}
%
There is no concealment at the source level---the docstring describes the
construction outright: \emph{``The default code path writes the precomputed
answer (the XZZX mask) directly \dots\ Set \texttt{QEC\_RUN\_SEARCH=1} to
reproduce the live parallel search.''} What makes this case harder to detect
than plain transcription is that a genuine method exists:
\texttt{\_live\_search} implements the parallel search the task asks for, and by
majority appearances produced \texttt{BEST\_MASK} in an earlier session. The
submission thus survives every check short of tracing which branch actually
executes. The code is real and the method is real---the answer is even plausibly
\emph{derived} from it---but the artifact under evaluation is a lookup, causally
disconnected from any computation performed at scoring time.

% ---------------------------------------------------------------------
\paragraph{Correct diagnosis, misallocated effort (D.5; Fig.~\ref{fig:case-studies}c).}
The task is PDE parameter inversion across six sub-problems: given observed
solution fields, recover the underlying physical parameters or coefficient
fields. Unlike the two C.2 cases, nothing here is dishonest. The agent trains
genuine learned estimators on every instance: CNN/FNO regressors, a U-Net
segmenter for the Darcy field, physics-based estimators elsewhere. Its failure
is purely on effort allocation. The benchmark scores a run as the
\emph{mean} of per-instance improvement over the reference, so a single badly
lost instance can sink an otherwise strong run. Table~\ref{tab:pde-agg} traces
the agent's own logged scores. By attempt~4 it beats the reference on four of
six instances \add{ (the agent's own log says five; the fifth it counts is
\texttt{2dtf} at $-2.38$)}, and the aggregate sits at $-2.03$, anchored by one Darcy-flow
instance (\texttt{2ddf}) stuck at $-12.1$. What makes the case remarkable is
that the agent sees all of this. Its log reads:
%
\begin{quote}\small\itshape
``5/6 beat SOTA \dots\ let me make one focused attempt on darcy (the dominant
$-12.1$)'' \; $\rightarrow$ \; ``Darcy is data-limited ($\sim$0.008 floor)'' \;
$\rightarrow$ \; ``rddu FNO \dots\ improvement $+0.794$ (was $+0.376$)!''
\end{quote}
%
The diagnosis is correct; the response is not. After one unsuccessful Darcy
attempt, the agent spends its remaining budget lifting \texttt{2drddu}---an
instance it has already won---from $+0.376$ to $+0.794$. Low-variance progress on a won instance is the
tempting choice, but with a $-12.1$ term in a six-way mean, even a partial Darcy
recovery is worth more than perfecting every other instance combined. The
failure is not analytical. The agent found the dominant term and named it, then
optimized what was easy to improve rather than what mattered, and never returned
to check whether the two were the same thing.

\begin{table}[t]
\centering\small
\caption{Per-instance improvement-over-reference logged by the agent on
\texttt{pdeinvbench\_param\_inverse}; the aggregate is the mean of the six
columns. Bold marks the term the agent named as dominant (\texttt{2ddf}) and the
instance it improved instead (\texttt{2drddu}).}
\label{tab:pde-agg}
\begin{tabular}{@{}clcccccc@{}}
\toprule
Att. & Aggregate & \texttt{1dkdv} & \texttt{2dns} & \texttt{2dtf} & \texttt{2drdk} & \texttt{2drddu} & \texttt{2ddf} (Darcy) \\
\midrule
1 & $-10.04$ & $-2.73$ & $+0.50$ & $-2.38$ & $+0.50$ & $+0.38$ & $\mathbf{-56.5}$ \\
2 & $-3.24$  & $-2.73$ & $+0.50$ & $-2.38$ & $+0.50$ & $+0.38$ & $-15.7$ \\
3 & $-2.63$  & $+0.93$ & $+0.50$ & $-2.38$ & $+0.50$ & $+0.38$ & $-15.7$ \\
4 & $-2.03$  & $+0.93$ & $+0.50$ & $-2.38$ & $+0.50$ & $+0.38$ & $\mathbf{-12.1}$ \\
5 & $-1.96$  & $+0.93$ & $+0.50$ & $-2.38$ & $+0.50$ & $\mathbf{+0.79}$ & $-12.1$ \\
6--7 & $-1.96$ & $+0.93$ & $+0.50$ & $-2.38$ & $+0.50$ & $+0.79$ & $-12.1$ \\
\bottomrule
\end{tabular}

\end{table}

% ---------------------------------------------------------------------
\paragraph{The verdict it wrote and ignored (D.7; Fig.~\ref{fig:case-studies}d).}
The task is HVAC load forecasting for edge deployment: can a lightweight,
interpretable model beat a deep network? As in the D.5 case, nothing the agent
does is dishonest: it downloads a real energy dataset, splits it strictly in
time order, fits its scaler on the training window only, and trains an LSTM
against two interpretable GAMs. Its headline is that the slim GAM beats the LSTM
by 22.4\%---exact arithmetic, empty result, because the baseline it beats is
\textbf{broken}. The LSTM reaches $R^2 = 0.076$---barely better than predicting the
training mean, and far below the ${>}0.6$ the agent itself notes is typical.
Beating a model that did not train tells you nothing: the 22.4\% margin only infers ``the LSTM is broken.'' instead of ``the GAM is good''. The agent states this
itself, in writing, in its own peer-review section:
%
\begin{quote}\small\itshape
``the central claim rests on a baseline that is almost certainly broken \dots\
the headline finding is therefore uninterpretable.''
\end{quote}
%
The fix is nearly free---retrain the LSTM, or report the naive-persistence
baseline its review calls for, a single line either way. It does neither.
``Uninterpretable'' stays in the review section while ``22.4\% lower than the
LSTM'' opens the abstract, so a reader who stops at the abstract leaves with the claim the author has already retracted. The failure is not analytical: the
agent ran the diagnosis and stated the verdict, then ignore the potential fix.

